\subsection{Авиаперелёты}

\subsubsection{Описание решения}

Во входных данных задачи даётся матрица весов, которой можно задать полноценный граф. Чтобы определить величину максимального ребра на минимальном пути между двумя городами \textit{(с точки зрения суммарного расстояния)} воспользуемся модификацией алгоритма Флойда-Уоршелла.

Заменим метрику суммарного расстояния величиной ребра с максимальным весом: это возможно поскольку обе операции ассоциативны, то есть могут быть применены в разной группировке одного пути. Алгоритм перебирает все доступные варианты путей, поэтому гарантированно минимизирует указанную метрику.

Ответ к задаче --- максимум в матрице посчитанных метрик.

\subsubsection{Асимптотическая сложность}

\textit{Временная сложность решения} --- кубическая \( \mathcal{O}(N^3) \). Алгоритм перебирает все возможные тройки вершин и считает метрики для полученных путей за константное время. Определение максимума в матрице происходит за квадратичное время.

\textit{Пространственная сложность решения} --- квадратичная \( \mathcal{O}(N^2) \). В алгоритме материализована матрица весов, которая переиспользуется для подсчёта матрицы метрик.

\subsubsection{Листинг кода}

\begin{lstlisting}
#include <algorithm>
#include <cstddef>
#include <iostream>
#include <vector>

class Graph {
public:
  explicit Graph(std::size_t n_count)
      : distances_(n_count + 1, std::vector<std::size_t>(n_count + 1, 0)) {
  }
  void AddEdge(std::size_t from_id, std::size_t to_id, std::size_t weight) {
    this->distances_[from_id][to_id] = weight;
  }
  std::size_t FindExpensiveEdge(std::size_t n_count) {
    for (std::size_t third_vertex = 1; third_vertex <= n_count; third_vertex++) {
      for (std::size_t first_vertex = 1; first_vertex <= n_count; first_vertex++) {
        for (std::size_t second_vertex = 1; second_vertex <= n_count; second_vertex++) {
          const std::size_t other_path = std::max(
              this->distances_[first_vertex][third_vertex],
              this->distances_[third_vertex][second_vertex]
          );
          this->distances_[first_vertex][second_vertex] =
              std::min(this->distances_[first_vertex][second_vertex], other_path);
        }
      }
    }
    std::size_t expensive_edge = 0;
    for (std::size_t first_vertex = 1; first_vertex <= n_count; first_vertex++) {
      for (std::size_t second_vertex = 1; second_vertex <= n_count; second_vertex++) {
        expensive_edge = std::max(expensive_edge, this->distances_[first_vertex][second_vertex]);
      }
    }
    return expensive_edge;
  }

private:
  std::vector<std::vector<std::size_t>> distances_;
};

int main() {
  std::size_t n_count = 0;
  std::cin >> n_count;
  Graph graph(n_count);
  for (std::size_t row_idx = 0; row_idx < n_count; row_idx++) {
    for (std::size_t col_idx = 0; col_idx < n_count; col_idx++) {
      std::size_t weight = 0;
      std::cin >> weight;
      graph.AddEdge(row_idx + 1, col_idx + 1, weight);
    }
  }
  std::cout << graph.FindExpensiveEdge(n_count);
  return 0;
}
\end{lstlisting}